% !TeX root = ../nc-quadric-surface-moduli.tex
% 2026-11-21 / v1.39
\section{構造対象の忘却と楕円型の有限・無限ファイバー}
\label{sec:unpointed-elliptic-fiber}

全ての Abel 圏を一つの代数スタックに集めることは仮定しない。
ここでは、$\HH$ の同型類から、その基礎 Abel 圏の exact 同値類への
集合論的な忘却写像を考える。固定 $\cC$ でのこの写像のファイバーは、
許容される $\OO$ の同型類を
$\operatorname{Aut}^{\mathrm{ex}}(\cC)$ で割った集合である。
固定 $\cC$ の上のスタックの \emph{2-fiber} はこの商ではなく、
許容 $\OO$ の同型群oidであることに注意する。

\begin{lemma}[許容構造対象の数値類と一意性]
\label{lem:forgetting-numerical-rigidity}
$(\cC,\OO)\in\HH$ を固定し、
定理~\cite[\texttt{thm:normal-helix}]{QuadricIntrinsic2026} の foundation によって
$K_0(\cC)\simeq K_0(Q)$、$Q=\PP^1\times\PP^1$ と同定する。
$\cC$ 内の Gabriel--Krull rank-one torsion-free exceptional 対象 $L$ は
ある $a,b\in\mathbb Z$ に対して
\[
 \operatorname{ch}[L]
 =1+aH_1+bH_2+abH_1H_2
 =\operatorname{ch}[\OO_Q(a,b)]
\]
を満たす。この数値類を指定したとき、該当する $L$ の同型類は
高々一つである。特に、固定 $\cC$ 上の許容構造対象の
同型類の集合は高々可算である。
\end{lemma}
\begin{proof}
命題~\cite[\texttt{prop:first-orbit-object}]{QuadricIntrinsic2026} の Euler 形式の計算と
$r=1$ という rank 条件を使う。
$\operatorname{ch}[L]=r+aH_1+bH_2+cH_1H_2$ と書くと
\[
 1=\chi(L,L)=r^2+2rc-2ab=1+2c-2ab.
\]
従って $c=ab$。上の $K_0$ は捩れを持たず、
この Chern 指標で整数類が一意に定まる。

同じ類を持つ二つの対象 $L,L'$ を取る。
$t=S^{-1}[2]$ と置く。もし $\Ext^2(L,L')\ne0$ なら
Serre 双対性により $\Hom(L',t^{-1}L)\ne0$ であり、
rank-one torsion-free 性から $L'\hookrightarrow t^{-1}L$ が得られる。
命題~\cite[\texttt{prop:first-orbit-object}]{QuadricIntrinsic2026} の証明で用いた AZ 豊富性と
自由被覆・負切断消滅を任意の coherent 対象に適用すると、
$n\gg0$ で
$\dim\Hom(\OO,t^n M)=\chi(\OO,t^nM)$ が成り立つ。
両辺に $\Hom(\OO,t^n-)$ を施すと、Riemann--Roch により
\[
 (2n+a+1)(2n+b+1)
 \le(2n+a-1)(2n+b-1)
\]
となり、矛盾する。従って $\Ext^2(L,L')=0$。
$\chi(L,L')=1$ から $\Hom(L,L')\ne0$ であり、同様に単射
$L\hookrightarrow L'$ を得る。余核は $K_0$ 類が零なので
Hilbert 多項式も零である。
$t$ の AZ 豊富性により非零 coherent tails 対象には
十分高い次数に非零切断が存在するから、余核は零である。
\end{proof}

滑らかな楕円曲線 $E$ と次数二の線束 $L_0,L_1,L_2$ に対し、
線束の類を加法的に $\ell_i=[L_i]\in\operatorname{Pic}^2(E)$ と書く。
差を
\[
 \eta=\ell_1-\ell_0,\qquad \delta=\ell_2-\ell_0
 \quad\in\operatorname{Pic}^0(E)
\]
とする。$(E,L_0,L_1,L_2)$ が VdB 許容であることは、
この滑らかな場合には $\eta\ne0,\eta\ne\delta,\delta\ne0$
と同値である。実際、相異なる次数二の線束による二つの
二重被覆 $E\to\PP^1$ は異なる被覆対合を持つので、
積写像は閉埋込みとなる。

\begin{lemma}[楕円四つ組の二つの反射]
\label{lem:elliptic-two-reflections}
許容楕円四つ組の helix の添字を一つ進める操作 $T$ は
\[
 T:(\eta,\delta)\longmapsto(\delta-\eta,\delta)
\]
を誘導する。また、
\[
 \eta\ne2\delta
 \tag{R-open}\label{eq:elliptic-ruling-open}
\]
の下で、命題~\ref{prop:ruling-involution-mutation} による
もう一方の ruling の foundation も cubic 型であり、
対応する四つ組は
\[
 R:(L_0,L_1,L_2)\longmapsto
 (L_0L_1L_2^{-1},L_2,L_1).
\]
従って
\[
 R:(\eta,\delta)\longmapsto(2\delta-\eta,\delta),
 \qquad RT:(\eta,\delta)\longmapsto(\eta+\delta,\delta).
\]
二つの操作は、それぞれ underlying Abel 圏の exact 同値類を保つ。
$R$ は $\OO$ を保ち、$T$ は一般に $\OO$ を変更する。
\end{lemma}
\begin{proof}
helix の関係
$L_{i+3}=L_i^{-1}L_{i+1}L_{i+2}$ により、
$\ell_{i+2}-\ell_i=\delta$ は $i$ に依存しない。
従って $(L_1,L_2,L_3)$ の差は
$(\delta-\eta,\delta)$ である。
添字の変更は同じ $\mathbb Z$-代数の reindexing なので、
基礎 tails Abel 圏を変えない。

後半は ruling 交換の二写像を楕円曲線の切断で計算する。
$V_i=H^0(E,L_i)$、$W_2=H^0(E,L_0L_1)$、
$W_3=H^0(E,L_0L_1L_2)$ とすると、
元の foundation の二写像は
\[
 V_1\otimes V_0\longrightarrow W_2,\qquad
 V_2\otimes W_2\longrightarrow W_3
\]
という切断の積である。
次数二の線束の評価列を $L_0L_1$ で捩ると、
\[
 \ker(V_2\otimes W_2\longrightarrow W_3)
 \simeq H^0(E,L_0L_1L_2^{-1})
\]
を得る。ここでは $\deg(L_0L_1L_2^{-1})=2$ による
$H^1$ 消滅を用いた。
一方、$L_1$ の評価列を $L_0L_1$ で捩ると、
$V_0\to V_1^\vee\otimes W_2$ の余核は
$H^0(E,L_0L_1^2)$ に同定される。
この同定では二次元空間の
$V_1^\vee\simeq V_1\otimes(\det V_1)^{-1}$ に由来する
一次元のスカラー因子を許す。
従って交換後の二写像は、
$(L_0L_1L_2^{-1},L_2,L_1)$ から得る
隣接二つの切断乗法と一致する。

新しい隣接次数二線束が相異なる条件は
$\eta\ne2\delta$ および $\eta\ne\delta$ であり、
残る非周期性条件は $\delta\ne0$ である。
ゆえに式~\eqref{eq:elliptic-ruling-open} の下で
新しい四つ組も VdB 許容である。
命題~\ref{prop:ruling-involution-mutation} と
定理~\cite[\texttt{thm:normal-helix}]{QuadricIntrinsic2026} による qgr の再構成により、
これは元の $(\cC,\OO)$ のもう一方の normalized foundation である。
差の計算は直接に
$\eta'=2\delta-\eta,\ \delta'=\delta$ を与える。
\end{proof}

\begin{theorem}[楕円型の全標識と忘却ファイバー]
\label{thm:elliptic-marking-lattice-fiber}
標数零の代数閉体上の滑らかな楕円許容四つ組
$U=(E,L_0,L_1,L_2)$ に対し、
\begin{equation}\label{eq:elliptic-nonresonance}
 \delta\ne0,\qquad \eta\notin\mathbb Z\delta
\end{equation}
を仮定する。$\cC=\qgr D(U)$ と置き、元の normalized foundation により
$K_0(\cC)$ を $K_0(Q)$ と同定する。
\begin{enumerate}
\item 各 $(a,b)\in\mathbb Z^2$ に対し、
$\operatorname{ch}L(a,b)=\operatorname{ch}\OO_Q(a,b)$ を持つ
rank-one torsion-free exceptional 対象 $L(a,b)$ が一意に存在する。
これらがそのような対象を尽くし、全ての $(\cC,L(a,b))$ が $\HH$ に属する。
\item 許容構造対象の同型類を $\operatorname{Aut}^{\mathrm{ex}}(\cC)$ で割った
忘却ファイバーは $\mathbb Z/{\sim}$ と同定される。
整数 $n$ は $L(n,0)$ に対応し、同値関係は次の有限な幾何学的判定で与えられる：
\begin{equation}\label{eq:elliptic-exact-fiber-relation}
 n\sim m
 \quad\Longleftrightarrow\quad
 \exists\epsilon\in\Aut(E,0),\quad
 \epsilon(\delta)=\delta,\quad
 \epsilon(\eta+n\delta)\in
 \{\eta+m\delta,\ 2\delta-\eta-m\delta\}.
\end{equation}
原点の選択はこの条件を変えない。
\item $\delta$ が非捩れなら、このファイバーは可算無限である。
$\delta$ が位数 $r$ の捩れ元なら、ファイバーは有限であり、
\eqref{eq:elliptic-exact-fiber-relation} を $\mathbb Z/r\mathbb Z$ に課した商である。
特にその濃度は $r$ 以下である。
\end{enumerate}
本定理は幾何点上の対象と同値類についての主張であり、
相対スタック射の離散性や étale 被覆性を結論しない。
\end{theorem}
\begin{proof}
まず対象を全て構成する。補題~\ref{lem:elliptic-two-reflections} の操作により
$U_n=(RT)^nU$ は具体的に
\[
 U_n=(E,L_0,L_1\otimes(L_2L_0^{-1})^{\otimes n},L_2)
\]
と書け、その差は $(\eta+n\delta,\delta)$ である。
仮定~\eqref{eq:elliptic-nonresonance} は、正負いずれの反復でも
$\eta+n\delta$ が $0,\delta,2\delta$ になることを排除する。
途中の $T$ による差 $\delta-\eta-n\delta$ に対しても同じである。
従って全ての四つ組が許容であり、各 $U_n$ から得る pointed 対象を
同じ underlying Abel 圏 $\cC$ へ移せる。

この構成の $K_0$ 上の効果を確認する。
元の foundation の数値類は
\[
 (\OO_Q,\OO_Q(1,0),\OO_Q(1,1),\OO_Q(2,1))
\]
であり、四つ先の項は $\OO_Q(2,2)$ に対応する。
$T$ は構造対象を第一 ruling へ移し、二つの座標軸を交換する。
$R$ の二つの評価列は、foundation を $(O,X,M,Z)$ と書くと
\[
 [Y]=2[M]-[Z],\qquad [Z']=2[M]-[X]
\]
を与える。従って $R$ は構造対象を固定し、二つの軸を交換する。
特に $RT$ による foundation の数値類は
\[
 (\OO_Q(1,0),\OO_Q(2,0),\OO_Q(2,1),\OO_Q(3,1))
\]
である。帰納的に、$U_n$ の構造対象は
$\operatorname{ch}\OO_Q(n,0)$ の対象 $L(n,0)$ となる。
逆操作でも同じ計算が成り立つ。
これらは定理~\cite[\texttt{thm:all-admissible-inverse}]{QuadricIntrinsic2026} により許容構造対象である。

元の cubic helix の適合二周期から来る根 $s$ を選ぶ。
この根は全 helix 項を二つずつ移す。
四項の基底と $s^2=t$ によれば、その $K_0$ 作用は
$\OO_Q(1,1)$ の tensor 作用に一致する。
ゆえに
\[
 L(a,b)=s^b L(a-b,0)
\]
は指定数値類を持つ許容構造対象である。
補題~\ref{lem:forgetting-numerical-rigidity} により、
全ての rank-one torsion-free exceptional 対象の数値類はこの形であり、
同じ類の対象は一意なので (1) が従う。
また自己同値 $s$ による対角移動を除くと、
全ての忘却ファイバーの元を $L(n,0)$ で代表できる。

次にその重複を決める。
系~\ref{cor:two-intrinsic-foundations} により、pointed exact 同値は
二つの内在的 foundation の非順序対を保つ。
$U_n$ の二つの foundation は $U_n,RU_n$ であり、
仮定により両者は geometric cubic 型である。
Van den Bergh の四つ組の復元
\cite[arXiv v3, Theorem 5.5.10]{VdB2011} を用いると、
$L(n,0)$ と $L(m,0)$ が pointed 同値となる必要十分条件は、
$U_n$ が $U_m$ または $RU_m$ と曲線同型を介して同型になることである。
ここでは ordered foundation から full helix を反 Serre 移動で復元する。
逆に四つ組の同型は添字零を保つ tails 同値を与える。

楕円曲線の自己同型を translation と $\epsilon\in\Aut(E,0)$ に分ける。
translation は $\operatorname{Pic}^0(E)$ 上の差には作用しないので、
必要条件はちょうど \eqref{eq:elliptic-exact-fiber-relation} である。
十分性も成り立つ。差が一致した後は、次数二線束 $L_0$ の差を
translation で合わせれば三つの線束が全て一致する。
この translation の存在は、代数閉体上の
$[2]:E\to E$ の全射性による。従って (2) が従う。

$\delta$ が非捩れなら、$\epsilon(\delta)=\delta$ となる
原点を保つ自己同型は恒等に限る。
非自明な $\epsilon$ に対して $\epsilon-1$ は非零 isogeny であり、
その核の元は捩れだからである。
従って同値関係は $n=m$、または
\[
 2\eta=(2-n-m)\delta
\]
だけである。第二式が可能でも、ある整数 $q$ による反射
$n\mapsto 2-q-n$ の二元以下の軌道を同一視するだけなので、商は可算無限である。
$\delta$ の位数が $r$ なら、まず $n\equiv m\pmod r$ が同値を与え、
残る条件はこの有限巡回群上の \eqref{eq:elliptic-exact-fiber-relation} となる。
これで (3) を得る。
\end{proof}

\begin{corollary}[一般楕円型における忘却の可算無限ファイバー]
\label{thm:generic-elliptic-unpointed-fiber}
滑らかな楕円許容四つ組で $\delta$ が非捩れ、
$\eta\notin\mathbb Z\delta$ なら、忘却の同型類ファイバーは可算無限である。
さらに $2\eta\notin\mathbb Z\delta$ なら、整数 $n$ による上記の代表に重複はない。
特に幾何学的一般点ではこの最後の条件を満たす。
\end{corollary}
\begin{proof}
前定理と、\eqref{eq:elliptic-exact-fiber-relation} の非捩れの場合の記述による。
幾何学的一般点では独立な Picard パラメータに対し
全ての等式 $r\delta=0$ と $2\eta=q\delta$ を避ける。
\end{proof}

\begin{corollary}[余分な楕円自己同型がない捩れの場合の濃度]
\label{cor:elliptic-torsion-fiber-cardinality}
前定理の仮定に加え $\Aut(E,0)=\{1,-1\}$ とし、
$\delta$ の位数を $r\ge2$ とする。
$2\eta\notin\langle\delta\rangle$ なら忘却ファイバーは正確に $r$ 個である。
$2\eta=q\delta$ なら、その個数は
\[
 \frac{r+\#\{n\in\mathbb Z/r\mathbb Z:2n=2-q\}}2
 =\begin{cases}
   (r+1)/2,&r\text{ が奇数},\\
   r/2+1,&r\text{ が偶数で }q\text{ が偶数},\\
   r/2,&r\text{ が偶数で }q\text{ が奇数}.
  \end{cases}
\]
余分な自己同型を持つ楕円曲線には、前定理の有限群を含む判定式を使う。
\end{corollary}
\begin{proof}
$r>2$ なら $\delta$ を固定する $\epsilon$ は恒等だけなので、
重複がある場合は反射 $n\mapsto2-q-n$ による商である。
$r=2$ には $\epsilon=-1$ も許されるが、
$2\delta=0$ を \eqref{eq:elliptic-exact-fiber-relation} に代入すると
恒等と同じ二つの条件しか生じない。
反射の軌道数は固定点数を用いる表示式となり、合同式 $2n=2-q$ を解けばよい。
\end{proof}

\begin{remark}[既存研究と残る問題]
上の定理は条件~\eqref{eq:elliptic-nonresonance} の下で全標識と
幾何点ファイバーの同値判定を与えるが、
自然な $\mathbb Z$-torsor 構造や相対ファイバーの代数性は与えない。
Kitamura--Okawa の2022年 Oberwolfach 報告\cite{OkawaKitamura2022OWR}は
cubic $\mathbb Z$-代数の無標識同値に $D_\infty$ が現れることと
rank-one torsion-free line bundle の $\mathbb Z^2$ 分類を述べる。
同報告で詳細論文は当時準備中とされているので、
本節の証明はその未照合の全分類定理に依存させず、
本稿の二つの foundation と楕円評価列から構成した。
\end{remark}
